\documentclass[11pt]{article}
\usepackage{manual_docs_style}
\externaldocument{build/environment_profiles}[environment_profiles.pdf]
\externaldocument{build/model_run_specs_schema}[model_run_specs_schema.pdf]

\begin{document}

\docTitle{Stage: Simulate}
\phantomsection\label{docstart:simulation_stage}
This stage converts planned Simulink runs into materialized run artifacts.


The main driver is \code{workflows/simulate/run_pending_sims.py}. It reads the planned topology from \code{model_run_specs.json}, checks the extracted model metadata, reconciles what already exists on disk, and launches MATLAB only for the runs that are still pending.

\section*{Invocation}


\begin{DocCode}
python workflows/simulate/run_pending_sims.py [--model MODEL] [--overwrite] [--rerun-failed] [--raise-on-errors] [--debug] [--run-id RUN_ID] [--runs-root RUNS_ROOT] [--matlab-cache-root MATLAB_CACHE_ROOT] [--profile-timing] [--compute-features [FEATURES1 ...]]
\end{DocCode}

\begin{itemize}
\item \code{--model}: If omitted, the script iterates over \hyperref[sec:allowed-tsenv-models]{all allowed \name{} models}.
\item \code{--overwrite}: mark all runs as \code{not_run} and materialize them again
\item \code{--rerun-failed}: mark only runs currently recorded as \code{failed} as \code{not_run} and materialize them again
\item \code{--raise-on-errors}: fail immediately instead of recording \code{failed} status
\item \code{--debug}: creates an additional folder called \code{<uuid>/debug} is populated with \code{input_signals.mat} and \code{raw_simulation.pickle}
\item \code{--run-id}: it will run (and overwrite) only the specific entry defined at \code{RUN_ID}, and  update \code{model_record.json}
\item \code{--runs-root}: optional base directory for run artifacts. When set, each model writes to \code{RUNS_ROOT/<MODEL>/runs/} instead of \code{models/simulink/<MODEL>/runs/}. This is intended for local NVMe scratch storage.
\item \code{--matlab-cache-root}: optional base directory for MATLAB/Simulink generated cache files. If omitted, MATLAB uses its normal cache and preference locations.
\item \code{--profile-timing}: writes per-phase timing records to \code{profile\_timing.jsonl} in the selected run root.
\item \code{--compute-features}: compute \code{features.json} with features specified. The feature names must match the names in \code{docs/overleaf_paper/src/appendix/appendix_environment.tex::model::Problem-specific features.}
\end{itemize}

The environment variable \code{SIMSCAPE_CUSTOM_PATH} must also be defined to the custom simscape blocks. For example: \code{SIMSCAPE_CUSTOM_PATH=/csem/divr/users/tbe/repo/tsENV/simscape_custom}.
\section*{Inputs and Source of Truth}

For each model under \code{models/simulink/<MODEL>/}, the script requires:
\begin{itemize}
\item \code{model_run_specs.json}: the planned run topology. See \hyperref[sec:model-run-specs-schema]{the \code{model_run_specs.json} schema}. It is the source of truth for:
\begin{itemize}
  \item baseline UUIDs and their \code{baseline_parameters}
  \item child intervention UUIDs and their single-parameter deltas
  \item time-zero baseline UUIDs implied by each child intervention
  \item \code{baseline_parameters_hash}, \code{parameter_hash}, and \code{time0_baseline_hash}
\end{itemize}

\item \code{simulink_model_original.mdl}
\item \code{experiment_config.json}: the source of truth for:
  \begin{itemize}
  \item \code{sampling_rate_hz}
  \item \code{end_time_input_s}
  \item the ordered list \code{observable_signals}
  \end{itemize}
\item \code{features.py}: Used for computing features on the raw signal.
\item \phantomsection\label{sec:model-record-json-schema}\code{runs/model_record.json}: the runtime inventory of all materialized or pending runs. It is a flat object keyed by run UUID\@. Each entry represents exactly one of:
  \begin{itemize}
  \item a baseline run
  \item an intervention child run
  \item a time-zero baseline run
  \end{itemize}


Schema:

\begin{DocCode}
{
  "<baseline_uuid>": {
    "class_internal": "no_interevention",
    "class_agent_facing_name": "no intervention",
    "run_type": "baseline",
    "parameters_hash": "<32-hex-hash>",
    "status": "not_run|success|failed",
    "timestamp": "..."
  },
  "<intervention_uuid>": {
    "class_internal": "simulation_specific_parameter_0",
    "class_agent_facing_name": "simulation_specific_parameter_0_facing_name",
    "run_type": "intervention",
    "parameters_hash": "<32-hex-hash>",
    "status": "not_run|success|failed",
    "timestamp": "..."
  },
  "<time0_baseline_uuid>": {
    "class_internal": "",
    "class_agent_facing_name": "",
    "run_type": "time0_baseline",
    "parameters_hash": "<32-hex-hash>",
    "status": "not_run|success|failed",
    "timestamp": "..."
  }
}
\end{DocCode}

The \code{timestamp} field records the wall-clock time when the runtime state is written in ISO-8601 format.
\end{itemize}



\section*{Internal Steps}
The code first validates the inputs and then follows  a reconcile-materialize pattern for \code{model_record.json}, then runs the simulations.
It can be called at the same time for multiple models.
The following applies to all three run kinds: baseline runs, intervention child runs, and time-zero baseline runs. Each is materialized independently and tracked independently in \code{model_record.json}.
\subsection*{Validation Between \code{metadata.json}, \code{experiment_config.json} and \code{description_levels.json}}
The stage validates the checks described \hyperref[sec:environment-validation-rules]{here}.
\subsection*{Reconciliation and Skip Flow}

The script first reconciles \code{model_run_specs.json} with \code{model_record.json} as follows:

\begin{enumerate}
\item If \code{model_run_specs.json} is newer than \code{model_record.json}, it first runs:
\begin{DocCode}
python refresh_model_run_spec_hashes.py models/simulink/<MODEL>
\end{DocCode}
This recomputes the hashes in \code{model_run_specs.json} before reconciliation continues.
\item It computes from \code{model_run_specs.json} the complete set of planned run UUIDs and hashes
	\item If a run is currently marked as \code{success}, its status is downgraded to \code{not_run} if either of the following conditions holds:
	  \begin{itemize}
	  \item the data artifact for its UUID is missing, that is, \code{runs/<UUID>/data.parquet} does not exist;
	  \item the stored runtime hash in \code{model_record.json} does not match the expected hash from \code{model_run_specs.json} for that run kind:
	    \begin{itemize}
	    \item baseline run: compare \code{model_record.json::<baseline_uuid>:"parameters_hash"} against \code{model_run_specs.json:<baseline_uuid>:"baseline_parameters_hash"}
	    \item intervention child run: compare \code{model_record.json::<child_uuid>:"parameters_hash"} against \code{model_run_specs.json:<baseline_uuid>:children:<child_uuid>:"parameter_hash"}
	    \item time-zero baseline run: compare \code{model_record.json::<time0_baseline_uuid>:"parameters_hash"} against \code{model_run_specs.json:<baseline_uuid>:children:<child_uuid>:"time0_baseline_hash"}
	    \end{itemize}
	  \end{itemize}
	\item If a run is set to \code{not_run}, then its folder \code{runs/<UUID>/} is deleted before re-materialization.
\item Any folder \code{runs/<UUID>/} that is not present in the reconciled \code{model_record.json} inventory is also deleted.
\end{enumerate}


If no baseline, intervention child, or time-zero baseline run remains in \code{not_run}, the stage rewrites \code{model_record.json} and exits without launching MATLAB.  

\subsection*{Simulation Flow}

After reconciliation, the stage executes every run marked as \code{not_run}. The runtime uses a per-model filesystem lock. This means
\begin{itemize}
\item one model can have at most one active writer at a time
\item different models can be simulated concurrently by different processes
\end{itemize}

The simulation pipeline consists of five main steps, all defined in \code{shared/simulation.py}: 
\begin{enumerate}
\item Simulink Simulation: \code{raw_simulation = simulate_recipe(recipe, observable_signals)}.
The \code{recipe} argument is a dictionary with:
  \begin{itemize}
  \item \code{baseline_parameters}: required; equal to \code{model_run_specs.json:<baseline_uuid>.baseline_parameters}
  \item \code{intervention_parameters}: optional; present for intervention child runs and time-zero baseline runs; equal to either \code{model_run_specs.json:<baseline_uuid>.children.<intervention_uuid>.parameters}
  \item \code{intervention_mode}: one of \code{none}, \code{at_intervention_time}, or \code{from_time_zero}
  \item \code{end_time_input_s}: required; the expected simulation end time for the run
  \end{itemize}

The function returns \code{raw_simulation}, a nested dictionary keyed by observable signal. Each signal entry contains internal keys \code{timestamp} and \code{data}. The \code{timestamp} arrays do not need to be aligned across signals.
Internally, the function works as follows:
\begin{enumerate}
\item it creates a copy of \code{simulink_model_original.mdl} at \code{models/simulink/<MODEL>/runs/<UUID>/simulink_model.mdl}
\item it uses \code{save_model_with_workspace_values(model_workspace)} to assign the simulation specific parameters to the copied model. \code{model_workspace} is a dict with entries that match \code{recipe.baseline_parameters}
Note that all \code{simulation_specific_parameters} are applied, otherwise it will result in an error. Additionally 
\begin{itemize}
  \item \code{end_time_input_s} is applied  with an inflated value of 1.20*\code{recipe.end_time_input_s}.
  \item \code{local_solver_step_size} is applied with a value equal to  1/\code{sampling_rate_hz}
\end{itemize}

\item it calls \code{sim_the_model.m}, a MATLAB script that invokes the internal Simulink engine with the copied \code{simulink_model.mdl}
\begin{DocCode}
  sim_the_model "TimeVaryingParameters" time_varying_parameters
\end{DocCode}
where \code{time\_varying\_parameters} is a dict with five keys:
\begin{itemize}
\item \code{identifier}: block paths in the working Simulink model
\item \code{key}: variable names associated with each block path
\item \code{time}: a two array sequences of timestamps for each parameter with 0 and end time (except for the intervention parameter, for it is \code{0,intervention_time})
\item \code{values}: sequences of values paired with each timestamp sequence (for all but the intervened parameter is repeated initial value)
\item \code{seen}: sequences initialized to zero alongside each parameter trajectory
\end{itemize}
For parameters that do not change during the run, \code{time} and \code{values} contain the initial value repeated at \code{0} and \code{recipe.end\_time\_input\_s}. For interventions, they contain the intervention timestamp and the updated value trajectory passed to MATLAB.
\item Only keep the signals in \code{observable_signals} for computing \code{raw_simulation}
\end{enumerate}
\item call \code{features.py::compute_problem_specific_features(raw_simulation[:end_time_input_s])}.
It returns a dictionary with problem-specific features that are useful for solving the problem.
For computing each of these features we use \code{features.py} in the corresponding model folder.
If a model has no \code{features.py}, the feature dict is \code{{}}.
The returned dict is saved as \code{features.json} in the run directory.
For missing feature values, scalar events are stored as JSON \code{null}, and list-valued features are stored as empty lists.
\item \code{df = resample_signal_dict(raw_simulation, signal_type, sampling_rate_hz)} validates that continuous raw signals match the target \code{sampling_rate_hz}, preserves impulse-like events with the configured event handling, and returns a dataframe.

Note that since the raw Simulink Simulation do not start from 0, the first available sample is at time \code{dt=1/sampling_rate_hz}.
Additionally we cut the simulation to \code{raw\_simulation[:end\_time\_input\_s]} with \code{end\_time\_input\_s} excluded.
\code{df} is expected to have \code{dt} as its first point and \code{floor(end\_time\_input\_s * sampling\_rate\_hz) - 1} points.

\item \code{serialize(df, uuid)} saves the dataframe in \code{runs/<UUID>/data.parquet}. 
All signal columns and the serialized \code{time} column are stored with \code{float32} precision.
If there is any \code{overflow} warning during the casting, this will be thrown as an error.
No casting to \code{float32} must be done anywhere else outside this function, including signal values, timestamps, target sample times, or resampling bin boundaries.


\item \code{validate(uuid)} validates the persisted artifact for that UUID after serialization by loading the saved dataframe from disk. Specifically, it checks that:
  \begin{itemize}
  \item the first timestamp available is \code{dt}
  \item columns are stored as \code{float32}
  \item the columns are named according to \code{experiment_config.json::observable_signals::observable_signals}, in the same order and the last column is time.
  \item the samples are regularly spaced according to \code{sampling_rate_hz}
  \item there are at least \code{50} points.
  \item there must be \code{floor(end_time_input_s * sampling_rate_hz)} -1 points
  \item the last point is smaller than \code{end_time_input_s}
  \item there is no inf, or nan value in the time series
  \end{itemize}

If any of these conditions are not met, the function raises an error.
\end{enumerate}
\section*{Output Artifacts}

By default the materialized artifacts are stored under:

\begin{DocCode}
models/simulink/<MODEL>/runs/<UUID>/
\end{DocCode}

When \code{--runs-root} is provided, artifacts are stored under:

\begin{DocCode}
<RUNS_ROOT>/<MODEL>/runs/<UUID>/
\end{DocCode}

Each successful run directory contains:
\begin{itemize}
\item \code{data.parquet}: the materialized time series
\item \code{simulink_model.mdl}: the exact Simulink model used to generate this run
\item \code{features.json}: the problem-specific feature dictionary computed from \code{observable_signals}. This file is omitted when \code{--compute-features} is not passed.
\end{itemize}

If \code{--debug} is passed, each successful run directory also contains:
\begin{itemize}
\item \code{debug/input_signals.mat}: the MATLAB debug payload storing \code{TimeVaryingParameters} (saved directly in \code{sim\_the\_model.m} before it is fed to the model)
\item \code{debug/raw_simulation.pickle}: the exact Python object returned by \code{simulate_recipe(...)}
\end{itemize}

\code{model_record.json} is rewritten under the run root and serves as the runtime inventory of completed, pending, or failed UUIDs. Any validation error or simulation error will lead to a failed UUID



\section*{Validation Rules}

In addition to the run-specific validation described above, the stage enforces the following operational rules:

\begin{itemize}
\item every materialized UUID must appear in the reconciled \code{model_record.json}
\item every \code{success} entry must have a corresponding \code{runs/<UUID>/data.parquet} artifact
\item every stored hash in \code{model_record.json} must match the hash expected from \code{model_run_specs.json} for that run kind
\item no stale run folder may remain under \code{runs/} after reconciliation
\end{itemize}


\end{document}
